\begin{center}
\hypertarget{native-062-article045}{}\pdfbookmark[2]{45.}{article-045}45.
\end{center}
Pars posterior tabulae primae operi huic annexae ad motum hyperbolicum spectat, ut iam supra monuimus, et pro argumento $A$ (utrique tabulae parti communi) logarithmum ipsius $B$ atque quantitatem $C$ ad septem figuras decimales (cifris praecedentibus omissis), quantitatem $T$ vero ad quinque dein ad sex figuras sistit. Extensa est haec pars, perinde ut prior, usque ad $A=0{,}300$, cui respondet $T=0{,}241207$, $u=2{,}930$ vel $=0{,}341$, $F=\pm52^0\,19'$; ulterior extensio superflua fuisset (art. 36).

Ecce iam ordinem calculi tum pro determinatione temporis ex anomalia vera tum pro determinatione anomaliae verae ex tempore. In problemate priori habebitur $T$ per formulam $T=\frac{e-1}{e+1}\operatorname{tang}\tfrac{1}{2}v^2$; ex $T$ tabula nostra dabit $\log B$ et $C$, unde erit $A=\frac{(1+C)T}{1-\tfrac{4}{5}T}$; hinc tandem per formulam 2) art. praec. invenietur $t$. In problemate posteriori computabuntur primo logarithmi constantium
\[
\begin{aligned}
\alpha&=\frac{75k\BareRoot{(\tfrac{1}{5}+\tfrac{9}{5}e)}}{2q^{\frac{3}{2}}}\mbox{\HistoricalDisplayNoteMark}\\[.7em]
\beta&=\frac{5e-5}{1+9e}\\[.7em]
\gamma&=\BareRoot{\frac{5e+5}{1+9e}}\mathord{\raisebox{.35ex}{.}}
\end{aligned}
\]
\noindent\HistoricalDisplayNoteText{Handschriftliche Bemerkung, welche sich auf eine von \textsc{Gauss} zu seinem eignen Gebrauch im Jahr 1816 berechnete noch nicht gedruckte Tafel bezieht, auf die in diesem Bande zurückgekommen werden wird:] in tabula nostra $\alpha=\frac{k\BareRoot{(\tfrac{1}{5}+\tfrac{9}{5}e)}}{q^{\frac{3}{2}}}\mathord{\raisebox{.35ex}{.}}$}
\SourceLeaf{063}Tunc determinabitur $A$ ex $t$ prorsus eodem modo ut in motu elliptico, ita scilicet ut motui medio $\frac{\alpha t}{B}$ in tabula \textsc{Barker}i respondeat anomalia vera $w$ atque fiat $A=\beta\operatorname{tang}\tfrac{1}{2}w^2$; eruetur scilicet primo valor approximatus ipsius $A$ neglecto vel si subsidia adsunt aestimato factore $B$; hinc tabula nostra suppeditabit valorem approximatum ipsius $B$, cum quo operatio repetetur; valor novus ipsius $B$ hoc modo prodiens vix unquam correctionem sensibilem passus, neque adeo nova calculi iteratio necessaria erit. Correcto valore ipsius $A$ e tabula desumetur $C$, quo facto habebitur
\[
\begin{aligned}
\operatorname{tang}\tfrac{1}{2}v&=\frac{\gamma\operatorname{tang}\tfrac{1}{2}w}{\BareRoot{(1+\tfrac{4}{5}A+C)}},\\[.7em]
r&=\frac{(1+\tfrac{4}{5}A+C)q}{(1-\tfrac{1}{5}A+C)\cos\tfrac{1}{2}v^2}\mathord{\raisebox{.35ex}{.}}
\end{aligned}
\]
Patet hinc, inter formulas pro motu elliptico et hyperbolico nullam omnino differentiam reperiri, si modo $\beta$, $A$ et $T$ in motu hyperbolico tamquam quantitates negativas tractemus.
